This study presented a comprehensive framework for UAV energy consumption modeling, leveraging a publicly available dataset from 209 autonomous flights of a DJI Matrice 100 quadcopter collected over seven months under diverse operational and environmental conditions. A threshold-based flight phase classification scheme was developed, using histograms of vertical velocity, horizontal speed, and power consumption to define distinct segments i.e. Taxi, Hover, Cruise, Climb, Descent, and Other, ensuring robust and physically meaningful segmentation of flight data.
A physics-based power model was formulated using actuator disk theory and phase-dependent efficiency estimates. The model demonstrated consistent performance across hover, climb, descent, and cruise phases, with MAPE values below 12\% in cruise flight, confirming its ability to capture fundamental aerodynamic and propulsion dynamics under steady-state conditions.
A Random Forest regression model was trained on 19 flight-state features including velocity, wind, attitude, and flight mode indicators and achieved an
$R^2$ score of $0.941$, with MAE consistently below 50~W across all flight phases. The model significantly outperformed the physics-based approach, reducing prediction errors by up to 50\% and accurately tracking transient power variations during complex maneuvers.
However, the data-driven model exhibited unphysical extrapolation behavior when exposed to out-of-distribution conditions, and its internal decision logic remained opaque, limiting interpretability. In contrast, the physics-based model, while less accurate in dynamic regimes, maintained physical consistency and robustness beyond the training domain.
These results confirmed a fundamental trade-off: data-driven models offer superior accuracy within their training domain, but at the cost of generalization and explainability. Physics-based models, though less precise in complex scenarios, provide greater transparency and reliability under uncertainty. Future work should focus on improving the representation of the hover phase in physical models. To enhance generalization in data-driven approaches, models must be trained on diverse datasets that encompass a broad range of parameter values and operational scenarios.